% ─────────────────────────────────────────────────────────────────────────────
% Supplementary: calibration protocol. Condensed from docs/04_calibration_protocol.md (full record, every number from there).
% ─────────────────────────────────────────────────────────────────────────────
\section{Calibration protocol}
\label{sec:calibration}

Apart from the attention layers, all hyperparameters of \method were set on two held-out calibration sets. Apart from choosing those layers, the test sets were used only to report results and to diagnose errors; any change suggested by a diagnosis had to pass on the calibration sets first.

\paragraph{Calibration sets.}
\emph{InfoVQA train} consists of 600 questions sampled with a fixed seed from the train split (ANLS).
\emph{MME-RW calib} is drawn from the full English MME-RealWorld (23,609 questions): we remove the questions of
the Lite subset (1,787) and all other questions on its images (2,654), and sample 120 of the remaining 19,168
questions from each of the five subtasks (accuracy, 600 questions). The result shares no questions and no images with the Lite test set.
DocVQA has no train split with answers, so InfoVQA train also represents documents.

\paragraph{Metrics.}
Retention is the score divided by the uncompressed score on the same set (83.9 and 59.5). We report its average
over the two sets and its worst value over six groups (InfoVQA and the five subtasks). TTFT is the median over 50
images after one warm-up image, with one process per GPU. Comparisons use a paired bootstrap on the same questions
with 2,000 resamples (4,000 in the candidate tests) at 95\% confidence.

\begin{table}[t]
  \centering
  \footnotesize
  \setlength{\tabcolsep}{4pt}
  \begin{tabular}{l ccccc c}
    \toprule
    $m$ & 1024 & 1536 & 2048 & 2560 & 3072 & Mean \\
    \midrule
    $1$ & 57.9 & 61.2 & 63.9 & 65.2 & 68.4 & 63.3 \\
    $K/2$ & 63.9 & 68.4 & 68.8 & 70.7 & 71.0 & \textbf{68.5} \\
    $3K/4$ & 65.2 & 67.6 & 68.3 & 70.0 & 70.9 & 68.4 \\
    $K$ & 63.5 & 68.8 & 69.1 & 70.1 & 71.8 & 68.6 \\
    \bottomrule
  \end{tabular}
  \caption{\textbf{Stage A: coverage fraction.} Score averaged over the two calibration sets, per budget $B$.}
  \label{tab:calib_m}
\end{table}

\begin{table}[t]
  \centering
  \footnotesize
  \setlength{\tabcolsep}{2pt}
  \begin{tabular}{l ccc ccc ccc}
    \toprule
    & \multicolumn{3}{c}{$m=K/2$} & \multicolumn{3}{c}{$m=3K/4$} & \multicolumn{3}{c}{$m=K$} \\
    \cmidrule(lr){2-4}\cmidrule(lr){5-7}\cmidrule(lr){8-10}
    $B$ & Avg. & Min & TTFT & Avg. & Min & TTFT & Avg. & Min & TTFT \\
    \midrule
    1024 & 89.7 & 76 & 0.90 & 91.2 & 77 & 0.69 & 88.5 & 80 & 0.40 \\
    1536 & 95.7 & 85 & 0.99 & 94.3 & 72 & 0.89 & 96.2 & 86 & 0.45 \\
    2048 & 95.9 & 85 & 1.08 & 95.5 & 72 & 1.00 & 96.0 & 74 & 0.56 \\
    2560 & 98.9 & 94 & 1.19 & 97.9 & 82 & 1.15 & 97.7 & 69 & 0.69 \\
    3072 & 99.3 & 86 & 1.41 & 99.1 & 83 & 1.32 & 100.1 & 78 & 0.80 \\
    \bottomrule
  \end{tabular}
  \caption{\textbf{Retention (\%) and TTFT (s)} for the three plausible coverage fractions, before the window was calibrated. Avg.: average over the two sets. Min: worst of the six groups.}
  \label{tab:calib_ret}
\end{table}

\paragraph{Stage A: coverage fraction and default budget.}
We swept $m\in\{1,K/4,K/2,3K/4,K\}$ with a provisional window. Using only the thumbnail ($m=1$) loses about 5
points at every budget (\cref{tab:calib_m}), while $m/K$ from 0.5 to 1 forms a plateau whose differences are mostly not significant.
Since the averages tie, we chose by stability: $m=K/2$ has the most stable worst-group retention (\cref{tab:calib_ret}). We set
$B=1536$ as the default because retention saturates there (90\% at $B=1024$, 96\% at $B=1536$, no further gain at $B=2048$).
Pure coverage ($m=K$) is about twice as fast and is reported as an ablation.

\begin{table*}[t]
  \centering
  \footnotesize
  \setlength{\tabcolsep}{4pt}
  \begin{tabular}{l c c c c}
    \toprule
    $B$ & $S$ & Avg.\ retention (\%) & Worst group (\%) & TTFT (s) \\
    \midrule
    1024 & 0 / \textbf{5} / 10 / 15 / 20 & 85.9 / \textbf{89.6} / 89.7 / 88.6 / 89.3 & 61.5 / \textbf{77.7} / 75.5 / 73.8 / 73.8 & 0.43 / \textbf{0.67} / 0.89 / 1.12 / 1.35 \\
    1536 & 0 / \textbf{4} / 9 / 13 / 18 & 93.2 / \textbf{95.6} / 95.7 / 95.5 / 94.5 & 80.0 / \textbf{85.1} / 84.6 / 89.2 / 78.5 & 0.58 / \textbf{0.77} / 0.99 / 1.18 / 1.41 \\
    2048 & 0 / \textbf{4} / 8 / 12 / 16 & 95.1 / \textbf{96.7} / 95.9 / 95.5 / 96.2 & 75.4 / \textbf{84.6} / 84.6 / 81.5 / 83.1 & 0.72 / \textbf{0.92} / 1.08 / 1.26 / 1.45 \\
    2560 & 0 / \textbf{3} / 7 / 10 / 14 & 97.7 / \textbf{98.8} / 98.9 / 97.5 / 97.5 & 89.2 / \textbf{93.6} / 93.8 / 83.1 / 80.0 & 0.86 / \textbf{1.02} / 1.19 / 1.30 / 1.48 \\
    3072 & 0 / \textbf{3} / 7 / 10 / 14 & 98.5 / \textbf{98.2} / 99.3 / 98.3 / 100.3 & 86.2 / \textbf{86.2} / 86.2 / 86.2 / 89.2 & 1.05 / \textbf{1.23} / 1.41 / 1.53 / 1.72 \\
    \bottomrule
  \end{tabular}
  \caption{\textbf{Stage B: scoring window} with $m=\lfloor K/2\rfloor$. The adopted value of $S$ is in bold.}
  \label{tab:calib_s}
\end{table*}

\paragraph{Stage B: scoring window.}
We fixed $m=\lfloor K/2\rfloor$ and swept $S\in\{0,S_0/2,S_0,1.5S_0,2S_0\}$, where $S_0$ is the provisional window
(\cref{tab:calib_s}). Removing the window ($S=0$) is significantly worse on local-detail questions, \eg $-$10.8
points [$-$18, $-$4] on Remote Sensing at $B=1024$ and $1536$. Larger windows bring no reliable gain and add TTFT.

\paragraph{Attention layers.}
The layers $\mathcal{L}_p=\{21,\dots,28\}$ and $l_q=20$ were chosen on InfoVQA val before this protocol existed, and
InfoVQA val is also a test set. They were not re-selected on the calibration sets. For the 4B and 14B models they are mapped by
relative depth, by a rule fixed before those models were run.

\begin{table}[t]
  \centering
  \footnotesize
  \setlength{\tabcolsep}{3pt}
  \begin{tabular}{>{\raggedright\arraybackslash}p{0.33\linewidth} >{\raggedright\arraybackslash}p{0.57\linewidth}}
    \toprule
    Candidate & Result on the calibration sets \\
    \midrule
    Window grows with the grid & No $f\in\{0.2,\dots,0.5\}$ passes all criteria; local-detail $\Delta$ at $f=0.2$: +0.67 [$-$0.44, +1.78] \\
    Question-stem query, top-patch aggregation & Worse on local detail for all eight variants, \eg $-$3.33 for the stem query \\
    Options sorted before scoring & Local detail $-$0.19 [$-$0.93, +0.46]; retention $-$0.56 at $B=1536$, misses the limit of $-$0.5 \\
    Attention-head selection, question-free correction & Recall +4.2 on a synthetic set, but retention $-$1.60, $-$0.83, $-$0.85; local detail $-$2.22 [$-$4.35, $-$0.09] \\
    \bottomrule
  \end{tabular}
  \caption{\textbf{Rejected candidates.} All were tested with the criteria below, and each decision used only the calibration sets.}
  \label{tab:calib_rej}
\end{table}

\paragraph{Rejected candidates.}
Diagnosing errors on the test sets suggested the four changes in \cref{tab:calib_rej}. Before each test we fixed three
criteria: average retention at most 0.5 points below the frozen configuration at every budget, a significant gain on the
local-detail questions of MME-RW (Remote Sensing, Monitoring, Autonomous Driving; 360 questions), and at most
0.3\,s of added TTFT at $B=1536$. None of the candidates passed its criteria, so we kept the configuration as frozen.

\paragraph{Deviations from the pre-registered rules.}
\begin{enumerate}
  \item The rule for Stage B selected $S=\{5,13,4,3,14\}$: the smallest $S$ within one point of the best average and
        worst-group retention. Each MME-RW group has only 120 questions, so the rule was driven by noise, and the
        differences to $S_0/2$ are not significant. We therefore used $S_0/2$ at every budget.
  \item The sweep of the growing window added $f=0.2$, since larger $f$ is very costly on MME-RW. The criteria did not change.
  \item Sorted options were added after the results of the other query variants were known. The criteria did not change, and the candidate was rejected.
  \item An earlier configuration had its window table and coverage fraction tuned on test data. We discarded it and
        recalibrated as described here.
\end{enumerate}
